
%% ============================================================
%%  CHƯƠNG 2: CƠ SỞ LÝ THUYẾT
%% ============================================================
\chapter{Cơ sở lý thuyết}
\section{Hệ thống IDS}
\label{sec:ids}

Hệ thống phát hiện xâm nhập (IDS) là một thành phần quan trọng trong hạ tầng an ninh mạng hiện đại, cho phép xác định các hoạt động độc hại trong các mạng doanh nghiệp, đám mây, công nghiệp và IoT. Như đã trình bày trong Mục~\ref{sec:reason}, bối cảnh mối đe dọa tiếp tục phát triển với các cuộc tấn công ngày càng tinh vi.

\subsection{Phân loại IDS}
IDS có thể được phân loại theo nhiều tiêu chí khác nhau:

\textbf{Theo phương pháp phát hiện:}
\begin{itemize}
    \item \textbf{Signature-based IDS}: Phát hiện dựa trên các mẫu tấn công đã biết (chữ ký). Hiệu quả cao với các tấn công cũ nhưng không phát hiện được tấn công mới (zero-day).
    \item \textbf{Anomaly-based IDS}: Phát hiện dựa trên sự khác biệt so với hành vi bình thường của hệ thống. Có thể phát hiện tấn công mới nhưng tỷ lệ dương tính giả cao.
\end{itemize}

\textbf{Theo nguồn dữ liệu:}
\begin{itemize}
    \item \textbf{Network-based IDS (NIDS)}: Phân tích lưu lượng mạng để phát hiện tấn công.
    \item \textbf{Host-based IDS (HIDS)}: Giám sát hoạt động trên một máy chủ cụ thể, bao gồm log hệ thống, tiến trình và tệp tin.
\end{itemize}

\subsection{IDS dựa trên học sâu}
Những tiến bộ gần đây trong ML và DL đã cải thiện đáng kể khả năng phát hiện của IDS thông qua việc tự động học các mẫu lưu lượng phức tạp từ dữ liệu mạng quy mô lớn. Khi xem xét IDS như một bài toán phân loại, vậy thì ta cần một mô hình có khả năng phân loại các kiểu lưu lượng mạng, với đầu vào là các đặc trưng của gói tin như header, IAT, flag... và đầu ra là số kiểu dữ liệu cần phân loại.

\subsubsection{Mạng CNN trong phát hiện xâm nhập}
Convolutional Neural Networks (CNN) có khả năng tự động học các đặc trưng không gian từ dữ liệu thô thông qua các lớp tích chập (convolutional layers). Trong bài toán IDS, lưu lượng mạng có thể được biểu diễn dưới dạng ma trận 2 chiều (ví dụ: time $\times$ features) và CNN có thể học các mẫu tấn công cục bộ trong không gian đặc trưng này. Các nghiên cứu gần đây cho thấy CNN đạt hiệu quả cao trong việc phân loại lưu lượng mạng và phát hiện các cuộc tấn công khác nhau.

\subsubsection{Mạng GRU trong phát hiện xâm nhập}
Gated Recurrent Units (GRU) là một biến thể của mạng nơ-ron hồi tiếp (RNN), được thiết kế để giải quyết vấn đề vanishing gradient thông qua cơ chế cổng (gate mechanism). Trong bối cảnh IDS, GRU có khả năng mô hình hóa các phụ thuộc thời gian trong chuỗi các gói tin mạng, cho phép phát hiện các hành vi tấn công có tính chất tuần tự như port scanning và brute force.

\subsubsection{Thách thức trong triển khai}
Thách thức chính trong việc triển khai IDS dựa trên học máy tập trung là yêu cầu về lượng lớn dữ liệu huấn luyện được gán nhãn từ các bên tham gia, điều thường khó đạt được trong thực tế do các ràng buộc về quyền riêng tư và quy định. Ngoài ra, các mẫu lưu lượng mạng thay đổi theo thời gian, đòi hỏi các vòng trao đổi dữ liệu phải diễn ra liên tục. Chi phí tính toán để huấn luyện các mô hình học sâu cũng là một vấn đề đối với các môi trường hạn chế tài nguyên. Những thách thức này thúc đẩy việc áp dụng Học liên kết (Federated Learning), cho phép huấn luyện mô hình cộng tác trong khi vẫn bảo vệ quyền riêng tư dữ liệu.

\section{Học liên kết (Federated Learning)}
\label{sec:fl}

\subsection{Học liên kết tập trung (Centralized Federated Learning)}
Trong Học liên kết tập trung (Centralized FL), một máy chủ khởi tạo và duy trì một mô hình toàn cục trong khi các clients phân tán cùng tối ưu hóa mô hình này bằng dữ liệu riêng của họ. Khung giải pháp được công nhận rộng rãi nhất trong bối cảnh này là FedAvg \cite{fedavg}. Trong FedAvg, mô hình toàn cục $f^0$ được khởi tạo bởi máy chủ và gửi đến tất cả clients trong tập $K$. Với mỗi vòng liên kết $t$, các clients đồng thời tối ưu hóa các tham số $w^t$ của $f^t$ sử dụng tập dữ liệu riêng $D_k, k\in K$, $w_k^t \leftarrow ClientUpdate(k, w_t)$ trước khi truyền các tham số mô hình cục bộ $w_k^t$ đến máy chủ. Máy chủ sau đó thực hiện tổng hợp có trọng số trên tập các tham số client để tính toán tham số mô hình toàn cục cho vòng tiếp theo $w^{t+1}$ theo Phương trình~\ref{fedavg}.

\begin{equation}
\label{fedavg}
w^{t+1} = \sum_{k \in K}\frac{|D_k|}{|D|}w_k^t
\end{equation}

trong đó $D_k$ là tập dữ liệu cục bộ của client $k$, $|D_k|$ là kích thước của nó, và $|D|$ là tổng kích thước mẫu của tất cả clients tham gia trong vòng liên kết hiện tại $t$. Máy chủ gửi mô hình toàn cục đã tính toán đến tất cả clients và lặp lại sang vòng tiếp theo.

Trong bối cảnh phát hiện xâm nhập, mỗi client quan sát các khối lượng mẫu tấn công khác nhau; do đó, sự phân phối dữ liệu không đồng nhất cả về số lượng và lớp phân loại là một thách thức cố hữu đối với học liên kết. Điều kiện này, được biết đến rộng rãi với tên gọi dữ liệu không độc lập và đồng nhất (Non-IID), gây ra hiện tượng ``client drift'' \cite{he2022fedssd}. Với dữ liệu Non-IID, mục tiêu tối ưu hóa cục bộ của mỗi client không thể tương thích với mục tiêu toàn cục do sự thiên lệch cục bộ của tập dữ liệu riêng ở mỗi phần tử client \cite{acar2021feddyn}, làm chậm tốc độ hội tụ và giảm hiệu suất mô hình toàn cục.

Để giải quyết thách thức Non-IID, nhiều thuật toán FL đã được đề xuất bao gồm các phương pháp thường biết đến nhất như FedProx \cite{li2020fedprox}, FedMLB \cite{kim2022fedmlb} và FedDyn \cite{acar2021feddyn}. Các phương pháp này điều chỉnh mục tiêu tối ưu hóa cục bộ để hạn chế client drift, cải thiện hiệu quả của học liên kết trong môi trường Non-IID.

Mặc dù có khả năng chống chịu cao với điều kiện Non-IID, FL tập trung yêu cầu các client truyền tham số mô hình, gây ra những hạn chế nghiêm trọng trong triển khai thực tế. Các nghiên cứu gần đây như Deep Leakage from Gradients (DLG) \cite{zhu2019dlg} và improved Deep Leakage from Gradients (iDLG) \cite{zhao2020idlg} đã thành công trong việc tái tạo mẫu huấn luyện từ gradient với độ chính xác đáng báo động, trong khi \cite{geiping2020inverting} đề xuất một cuộc tấn công tái tạo mạnh hơn thiết kế trực tiếp cho môi trường FL, cho thấy rằng các tham số mô hình không còn bảo đảm quyền riêng tư trước một máy chủ không đáng tin. Ngoài tấn công đảo ngược mô hình, nhiều công trình đề xuất tấn công suy diễn \cite{pyrgelis2017knockknock, melis2019mia, ganju2018inference} trên các tham số mô hình, khai thác được các thông tin cá nhân nhạy cảm của từng phân tử tham gia.

Nhiều công trình đã tìm cách giải quyết điểm yếu này: SecureAggregation \cite{bonawitz2016secagg} và DeepFed \cite{li2020deepfed} là các biến thể bảo vệ quyền riêng tư dựa trên các kỹ thuật mã hóa như Diffie-Hellman masking hoặc mã hóa đồng cấu Pallier. Tuy nhiên, các phương pháp này yêu cầu chi phí tính toán cao đáng kể và các hạn chế trong môi trường production. Ví dụ, SecureAggregation tạo thêm các vòng truyền thông và cơ chế phục hồi dropout, làm tăng độ phức tạp giao thức trong triển khai quy mô lớn. Tương tự, các phương pháp dựa trên mã hóa đồng cấu như DeepFed tạo ra chi phí tính toán cao và lệ thuộc vào việc quản lý khóa giải mã.

\subsection{Chắt lọc kiến thức liên kết tập trung (Centralized Federated Distillation)}
\label{sec:fd}
Chắt lọc kiến thức liên kết (Federated Distillation -- FD) \cite{li2019fedmd, tan2022fedproto, jeong2018federateddistillation} nổi lên như một biến thể của Học liên kết kết hợp FL với Chắt lọc kiến thức (Knowledge Distillation) \cite{hinton2015kd}. Các giải pháp FD, bao gồm FedMD \cite{li2019fedmd} và FederatedDistillation \cite{jeong2018federateddistillation} truyền tải các biểu diễn tri thức trung gian như logits hoặc class prototypes \cite{tan2022fedproto} thay vì tham số mô hình.

FederatedDistillation tính toán vector logit trung bình đã kích hoạt softmax $\bar{p}_{c,k}$ của mỗi lớp $c$ trong tập dữ liệu riêng $D_k$ của mỗi client $k$, $\bar{p}_{c,k} \leftarrow \frac{1}{|c|}\sum_{x \in D_k, y|x = c}p(x)$. Sau đó, máy chủ tính tổng các vector logit của tất cả clients cho mỗi lớp $c \in C$, $\bar{p}_c \leftarrow \sum_{k \in K}\bar{p}_{c,k}$. Với mỗi lớp, mỗi client trừ vector trung bình của chính nó khỏi vector toàn cục để thu được vector trung bình leave-one-out $\hat{p}_{c,k} \leftarrow \frac{1}{|K|-1} (\bar{p}_c - \bar{p}_{c,k})$. Mỗi client tối ưu hóa mô hình cục bộ qua Phương trình \ref{federateddistillation} cho mỗi mẫu $x$ trong tập dữ liệu riêng:

\begin{equation}
\label{federateddistillation}
w_k \leftarrow w_k - \eta \nabla\{\phi(p(x),y_i) + \gamma \cdot \phi(p(x), \hat{p}_{c,k}\}
\end{equation}

trong đó $\phi$ là hàm mất mát cross-entropy, $\gamma$ là siêu tham số điều khiển cường độ chắt lọc.

Các khung giải pháp FD gần đây đã phát triển để hoạt động trên một tập dữ liệu công khai $D_p$, có thể truy cập bởi tất cả các thành viên liên kết. FedMD \cite{li2019fedmd} là một công trình phổ biến sử dụng tập dữ liệu công khai có nhãn. Trong cách tiếp cận của họ, tất cả clients huấn luyện mô hình riêng trên tập dữ liệu công khai cho đến khi hội tụ, sau đó chuyển sang tập dữ liệu riêng trước quá trình liên kết. Với mỗi vòng liên kết, clients tính toán logits chưa qua hàm softmax cho mỗi mẫu của tập dữ liệu công khai. Máy chủ tổng hợp các logits của clients để tạo thành một consensus tổng kết biểu diễn của các mô hình tham gia lên từng mẫu công khai. Mỗi client huấn luyện mô hình cục bộ để phân phối dự đoán nhãn công khai của nó khớp với consensus này. Trong giai đoạn Digest của FedMD, consensus của máy chủ đóng vai trò là mô hình teacher, trong khi mỗi mô hình cục bộ đóng vai trò là mô hình student. Mỗi mô hình student học cách khớp đầu ra của nó trên các mẫu công khai với mô hình teacher là consensus thông qua một hạng tử trong hàm mất mát.

SSFL-IDS \cite{zhao2022ssfl} tận dụng các nguyên lý của FD và học bán giám sát (semi-supervised learning). SSFL-IDS không yêu cầu tập dữ liệu công khai phải được gán nhãn. Trong cách tiếp cận của họ, mỗi client tối ưu hóa mô hình cục bộ bằng tập dữ liệu riêng và tính toán logits đã kích hoạt softmax trên các mẫu công khai. SSFL-IDS triển khai một mô hình phân loại nhị phân (discrminator) để loại bỏ các dự đoán trên các mẫu công khai mà mô hình chính không chắc chắn. SSFL-IDS dựa vào softmax score của mô hình chính trên dữ liệu cục bộ và dữ liệu công khai để chọn ra tập dữ liệu phân loại nhị phân, giúp mô hình discriminator phân biệt các dự đoán của mô hình chính lên các mẫu công khai thuộc nhãn lạ, qua đó mỗi client trong SSFL-IDS chỉ gửi các dự đoán mà mô hình có độ tự tin cao để server tổng hợp, giảm thiểu đặc tính lệch nhãn của môi trường Non-IID. Sau đó, phân phối softmax score của mỗi mẫu dữ liệu thông qua cơ chế majority voting để chọn ra một nhãn giả (pseudo label) cho mẫu công khai tương ứng. Khi đó, bộ dữ liệu công khai, với nhãn giả, đã trở thành một bộ dữ liệu học giám sát đầy đủ, và mỗi client sẽ huấn luyện mô hình trên tập dữ liệu gán nhãn giả này như là một kiểu học liên kết.

\subsection{Học liên kết phân tán (Decentralized Federated Learning)}
\label{sec:dfl}
Học liên kết phân tán (Decentralized Federated Learning -- DFL) loại bỏ sự phụ thuộc vào máy chủ trung tâm vĩnh viễn bằng cách phân phối trách nhiệm điều phối giữa các clients tham gia \cite{beltran2023dflsurvey}. Các iải pháp DFL hiện có được chia thành nhiều cấu trúc liên kết khác nhau, bao gồm kiến trúc vòng (ring), lưới (mesh), cụm (cluster) và hình sao (star). Các cấu trúc này khác nhau về độ phức tạp truyền thông, khả năng ổn định khi client dropout, tốc độ hội tụ của mô hình liên kết và thuật toán tổng hợp.

Trong số các cấu trúc này, DFL dạng hình sao duy trì hệ thống liên kết gần giống với FL tập trung nhất, nơi một điều phối viên tạm thời đứng ra thực hiện việc tổng hợp trong mỗi vòng liên kết. Không giống như FL tập trung, vai trò điều phối viên được chỉ định mới liên tục trong quá trình liên kết. BrainTorrent \cite{braintorrent} là một công trình tiêu biểu trình bày một giải pháp FedAvg phân tán cho kiểu giao tiếp ngang hàng (peer-to-peer) theo cấu trúc hình sao.

Trong bối cảnh phát hiện xâm nhập, các nghiên cứu gần đây đã khám phá DFL-based IDS. 2DF-IDS \cite{friha20232dfl-ids} đề xuất một IDS dựa trên FL phân tán cho IIoT với cơ chế tổng hợp hoàn toàn phân tán. DOF-IDS \cite{nakip2023dfl-ids} cho phép FL phân tán cho các nút IoT và IP với cách tiếp cận dựa trên đăng ký. Tuy nhiên, mặc dù các khung giải pháp này loại bỏ máy chủ trung tâm vĩnh viễn, hầu hết vẫn tiếp tục trao đổi tham số mô hình giữa các thành viên, để lại chi phí truyền thông và rủi ro đảo ngược mô hình như những thách thức đáng kể.

\section{Tấn công độc hại (Poisoning Attacks) và Phòng thủ}
\label{sec:poisoning}

Bản chất cộng tác của học liên kết khiến quá trình huấn luyện dễ bị tấn công bởi các thành viên độc hại làm hỏng quá trình tổng hợp. Thông qua việc thao túng dữ liệu huấn luyện cục bộ, bản cập nhật mô hình và không gian dự đoán, các cuộc tấn công đầu độc có thể làm suy giảm đáng kể hiệu suất mô hình toàn cục và đặt ra thách thức lớn cho việc triển khai thực tế các hệ thống phát hiện xâm nhập liên kết.

\subsection{Tấn công đầu độc dữ liệu (Data Poisoning)}
Tấn công đầu độc dữ liệu là một cuộc tấn công có thể được thực hiện trong tất cả các khung giải pháp liên kết. Nó thao túng tập dữ liệu huấn luyện cục bộ của các client Byzantine trước khi huấn luyện, thúc đẩy quá trình tối ưu hóa cục bộ tạo ra một mô hình bị đầu độc. Một ví dụ phổ biến trong bối cảnh phát hiện xâm nhập là đảo nhãn có chủ đích (targeted label flipping), nơi các mẫu từ lớp tấn công thường xuyên xuất hiện được gán nhãn lại thành lưu lượng lành tính, buộc mô hình học các ranh giới quyết định sai. Các client bị đầu độc sau đó gửi các tham số mô hình bị đầu độc hoặc tạo ra các biểu diễn sai lệch ảnh hưởng xấu đến mô hình tổng hợp hoặc consensus.

\subsection{Tấn công đầu độc mô hình (Model Poisoning)}
Tấn công đầu độc mô hình sửa đổi các bản cập nhật mô hình cục bộ trước khi tổng hợp toàn cục, mà không cần bất kỳ dữ liệu huấn luyện nào. Kẻ tấn công có thể tạo ra các bản cập nhật tùy ý làm hỏng nghiêm trọng mô hình được tổng hợp một cách ngây thơ, khiến nó hoàn toàn không đáng tin cậy và làm đình trệ quá trình liên kết.

PoisonedFL \cite{xie2025poisonedfl} là một cuộc tấn công đầu độc mô hình gần đây có khả năng thích ứng với các phòng thủ hiện đại, bao gồm FLTrust \cite{cao2020fltrust} và FLAME \cite{nguyen2022flame}. Tất cả các client Byzantine trong cuộc tấn công này tạo ra một bản cập nhật bị đầu độc giống hệt nhau để giảm hiệu ứng ``tự triệt tiêu'' (self cancellation), đồng thời thích ứng bằng cách giảm cường độ tấn công nếu phát hiện sự hiện diện của cơ chế phòng thủ.

\subsection{Tấn công đầu độc Logit (Logit Poisoning)}
Trong các hệ thống chắt lọc kiến thức liên kết, các client Byzantine trực tiếp thao túng các bản cập nhật biểu diễn ảnh hưởng xấu đến consensus tổng hợp. Tấn công tối đa hóa mất mát (Loss Maximization Attack -- LMA) được giới thiệu bởi \cite{roux2025byzantine} tạo ra các logits đối nghịch nhằm tối đa hóa sự bất đồng trong quá trình tổng hợp, làm hỏng mô hình teacher consensus, và từ đó làm hỏng các mô hình student vốn là các mô hình lành tính.

Các cuộc tấn công đầu độc gần đây đã tiến hóa từ việc làm suy giảm FL bằng các giá trị ngoại lai cực đoan rõ ràng đến việc tạo ra các bản cập nhật độc hại ẩn danh nằm gần với hành vi Non-IID lành tính. Kết quả là, việc phân biệt các thành viên độc hại khỏi các client không đồng nhất tự nhiên ngày càng trở nên khó khăn.

\subsection{Cơ chế phòng thủ (Anti-Poisoning)}
\label{sec:defense}

Để giảm thiểu rủi ro từ các cuộc tấn công Byzantine, nhiều công trình nghiên cứu FL đã đề xuất các cơ chế khác nhau nhằm giảm hoặc loại bỏ hoàn toàn ảnh hưởng độc hại của các bản cập nhật trong quá trình tổng hợp.

\textbf{Robust Aggregation Rule (RAR) đơn giản:}
\begin{itemize}
    \item \textbf{Coordinate-wise Median} \cite{yin2018fedcomed}: Chọn giá trị trung vị cho mỗi tham số, loại bỏ các giá trị ngoại lai cực đoan.
    \item \textbf{TrimmedMean} \cite{yin2018fedcomed}: Loại bỏ phần cao nhất và thấp nhất của các bản cập nhật client cho mỗi tham số.
    \item \textbf{Krum} \cite{blanchard2017krum}: Chọn bộ tham số gần nhất với tất cả các bộ tham số lân cận của nó.
    \item \textbf{Bulyan} \cite{guerraoui2018bulyan}: Kết hợp Multi-Krum và TrimmedMean trong một quy trình tổng hợp thành hai bước.
\end{itemize}

\textbf{Các phương pháp hiện đại:}
\begin{itemize}
    \item \textbf{FLAME} \cite{nguyen2022flame}: Máy chủ tính pair-wise cosine distance cho mỗi bản cập nhật client, phân cụm chúng bằng HDBSCAN \cite{mcinnes2017hdbscan}, loại bỏ các cụm ngoại lai, cắt tỉa các mô hình (clipping) và áp dụng nhiễu (noising).
    \item \textbf{Cronus} \cite{chang2019cronus}: Triển khai RobustFilter như một bộ tổng hợp logit chống chịu Byzantine, lọc các dự đoán ngoại lai trước khi gán nhãn giả (pseudo-labeling).
    \item \textbf{ExpGuard} \cite{roux2025byzantine}: Đo khoảng cách của mỗi bản cập nhật client so với kết quả tổng hợp của thuật toán tổng hợp RAR để tính điểm ngoại lai $\sigma_k$ cho mỗi client, và tổng hợp mô hình sau cùng với điểm ngoại lai là trọng số tổng hợp.
\end{itemize}

\subsection{Chắt lọc kiến thức dựa trên bằng chứng (Evidential Knowledge Distillation)}
\label{sec:ekd}

Chắt lọc kiến thức truyền thống truyền thụ tri thức thông qua các phân phối xác suất thu được bằng chuẩn hóa softmax. Mặc dù hiệu quả, softmax loại bỏ thông tin chứa trong độ lớn tuyệt đối của các logits, khiến các vector dự đoán với mức độ tin cậy khác nhau tạo ra các phân phối xác suất giống hệt nhau sau chuẩn hóa.

Chắt lọc kiến thức dựa trên bằng chứng (Evidential Knowledge Distillation -- EKD), được đề xuất bởi \cite{xiang2025ekd}, truyền độ bất định của mô hình teacher cùng với kiến thức của nó đến các mô hình student. Trong EKD, các raw logits được kích hoạt bằng hàm mũ $e=\text{exp}(z)$, kết hợp với trọng số tiên nghiệm $\lambda$ để tạo ra vector tham số Dirichlet $|C|$ chiều $\alpha = e + \lambda$. Mô hình student chắt lọc kiến thức từ mô hình teacher bằng cách tối ưu hóa:

\begin{equation}
\label{ekd}
\begin{split}
L & = L_{1st} + \gamma L_{2nd} \\
  & = KL\left(\frac{\alpha_i^T}{\sum_{i \in C}\alpha_i^T} \middle|\middle| \frac{\alpha_i^S}{\sum_{i \in C}\alpha_i^S}\right) \\
  & +\gamma \cdot KL(Dir(p^T|\alpha^T)||Dir(p^S|\alpha^S))
\end{split}
\end{equation}

\subsection{Phát hiện ngoài phân phối dựa trên năng lượng (Energy-based OOD Detection)}
\label{sec:energy}

Ước lượng độ tin cậy rất quan trọng trong việc xác định các dự đoán thiếu độ tự tin của mô hình. Các phương pháp truyền thống dựa trên xác suất softmax tối đa thường chịu hạn chế là kết quả ước lượng quá tự tin ngay cả đối với các mẫu ngoài phân phối \cite{liu2020energy}. Phương pháp phát hiện OOD dựa trên năng lượng được đề xuất bởi \cite{liu2020energy} định lượng độ bất định của mô hình trên một mẫu dữ liệu $x$ dưới dạng năng lượng:

\begin{equation}
\label{energy}
E(x) = -T \log\sum_i e^{{z_i}/T}
\end{equation}

Năng lượng thấp cho thấy dự đoán tin cậy do phân phối logit có độ lớn tập trung vào một kênh, trong khi giá trị năng lượng cao thường cho thấy một mẫu không chắc chắn hoặc ngoài phân phối do phân phối logit có độ lớn nhỏ và rải rác giữa các kênh.

\pagebreak
